\documentclass[11pt]{article}
\usepackage[utf8]{inputenc}
\usepackage[margin=1in]{geometry}
\usepackage{amsmath, amssymb, amsthm}
\usepackage{enumitem}
\usepackage{xcolor}
\usepackage[most]{tcolorbox}
\usepackage{hyperref}

% Define colored boxes for different sections
\newtcolorbox{problemconfigbox}{
    colback=blue!5!white,
    colframe=blue!75!black,
    title=\textbf{1. PROBLEM CONFIGURATION ANALYSIS},
    fonttitle=\bfseries,
    breakable
}

\newtcolorbox{strategicbox}{
    colback=pink!5!white,
    colframe=pink!75!black,
    title=\textbf{2. STRATEGIC FRAMEWORK APPLICATION},
    fonttitle=\bfseries,
    breakable
}

\newtcolorbox{tacticalbox}{
    colback=green!5!white,
    colframe=green!75!black,
    title=\textbf{3. TACTICAL METHODOLOGY SELECTION},
    fonttitle=\bfseries,
    breakable
}

\newtcolorbox{conversionbox}{
    colback=yellow!5!white,
    colframe=yellow!75!black,
    title=\textbf{4. CONVERSION TECHNIQUES APPLICATION},
    fonttitle=\bfseries,
    breakable
}

\newtcolorbox{verificationbox}{
    colback=red!5!white,
    colframe=red!75!black,
    title=\textbf{5. VERIFICATION STRATEGY DESIGN},
    fonttitle=\bfseries,
    breakable
}

\newtcolorbox{roadmapbox}{
    colback=cyan!5!white,
    colframe=cyan!75!black,
    title=\textbf{6. SOLUTION ROADMAP},
    fonttitle=\bfseries,
    breakable
}

\newtcolorbox{competitionbox}{
    colback=purple!5!white,
    colframe=purple!75!black,
    title=\textbf{7. COMPETITION STRATEGY INTEGRATION},
    fonttitle=\bfseries,
    breakable
}

\newtcolorbox{keyinsightbox}{
    colback=orange!10!white,
    colframe=orange!75!black,
    fonttitle=\bfseries,
    breakable
}

\newtcolorbox{warningbox}{
    colback=red!10!white,
    colframe=red!75!black,
    fonttitle=\bfseries,
    breakable
}

\newtcolorbox{tipbox}{
    colback=teal!10!white,
    colframe=teal!75!black,
    fonttitle=\bfseries,
    breakable
}

\title{\textbf{Strategic Analysis: AMC 12B 2016 Problem 19}}
\author{Comprehensive Framework Analysis}
\date{}

\begin{document}

\maketitle

\section*{Problem Statement}

Tom, Dick, and Harry are playing a game. Starting at the same time, each of them flips a fair coin repeatedly until he gets his first head, at which point he stops. What is the probability that all three flip their coins the same number of times?

\[\textbf{(A)}\ \frac{1}{8} \qquad \textbf{(B)}\ \frac{1}{7} \qquad \textbf{(C)}\ \frac{1}{6} \qquad \textbf{(D)}\ \frac{1}{4} \qquad \textbf{(E)}\ \frac{1}{3}\]

\vspace{1em}

\begin{problemconfigbox}
\subsection*{A. Problem Classification}
\begin{itemize}[leftmargin=*]
    \item \textbf{Problem Type}: To Find - compute the probability of a specific event
    \item \textbf{Difficulty Band}: Late-band (Problem 19 of 25)
    \item \textbf{Competition Context}: AMC12 (75 min, 25 problems) - approximately 3 minutes per problem
    \item \textbf{Mathematical Domain}: Probability Theory with elements of Geometric Series and Infinite Series
\end{itemize}

\subsection*{B. Problem Structure Identification}
\begin{itemize}[leftmargin=*]
    \item \textbf{Given Information}: 
    \begin{itemize}
        \item Three independent players (Tom, Dick, Harry)
        \item Each flips a fair coin repeatedly ($P(\text{H}) = P(\text{T}) = \frac{1}{2}$)
        \item Each stops when they get their first head
        \item All start simultaneously
    \end{itemize}
    \item \textbf{Constraints}: 
    \begin{itemize}
        \item Each player's sequence is independent
        \item Players stop after first success (geometric distribution)
        \item Number of flips is unbounded (could theoretically be infinite)
    \end{itemize}
    \item \textbf{Objective}: Find $P(\text{all three flip the same number of times})$
    \item \textbf{Hidden Assumptions}: 
    \begin{itemize}
        \item The coin flips are independent across players
        \item The coins are fair and identically distributed
        \item We need to consider all possible numbers of flips (from 1 to $\infty$)
    \end{itemize}
    \item \textbf{Answer Choice Leverage}: 
    \begin{itemize}
        \item All answers are simple fractions with small denominators
        \item Answer (A) $\frac{1}{8}$ is suspiciously simple (might be a trap - probability on first flip only)
        \item Answer (B) $\frac{1}{7}$ is unusual and might come from infinite series
        \item Can use answer choices to verify if computed series sum is reasonable
    \end{itemize}
\end{itemize}

\subsection*{C. Strategic Orientation}
\begin{itemize}[leftmargin=*]
    \item \textbf{Hypothesis}: Each player follows a geometric distribution with success probability $p = \frac{1}{2}$
    \item \textbf{Conclusion}: We need to sum probabilities over all possible flip counts where all three succeed on the same flip
    \item \textbf{Notation}: 
    \begin{itemize}
        \item Let $P_n$ = probability that a single player gets their first head on flip $n$
        \item Let $P(\text{all same})$ = target probability
        \item $n$ = number of flips (ranges from 1 to $\infty$)
    \end{itemize}
    \item \textbf{Intuitive Approach}: This is a sum over all possible outcomes where all three succeed simultaneously, suggesting an infinite geometric series
    \item \textbf{Cross-Domain Potential}: 
    \begin{itemize}
        \item Geometric probability
        \item Infinite series and convergence
        \item Could use recursive/self-similarity approach
        \item Generating functions (advanced)
    \end{itemize}
\end{itemize}
\end{problemconfigbox}

\begin{strategicbox}
\subsection*{A. Psychological Strategies}
\begin{itemize}[leftmargin=*]
    \item \textbf{Mental Toughness}: This is a late-band problem requiring comfort with infinite series; don't be intimidated by the infinite sum - it telescopes nicely
    \item \textbf{Creativity Cultivation}: Consider both direct infinite series approach AND recursive self-similarity approach (two elegant solutions available)
    \item \textbf{Iterative Refinement}: If direct computation seems messy, step back and look for self-similar structure in the problem
\end{itemize}

\subsection*{B. Investigation Strategies}
\begin{itemize}[leftmargin=*]
    \item \textbf{Startup Orientation}: Recognize this as a geometric probability problem; identify that we need to sum over all possible synchronized outcomes
    \item \textbf{Get Your Hands Dirty}: 
    \begin{itemize}
        \item Compute probability for $n=1$: All three get heads on first flip = $\left(\frac{1}{2}\right)^3 = \frac{1}{8}$
        \item Compute for $n=2$: All three get T then H = $\left(\frac{1}{2}\right)^6 = \frac{1}{64}$
        \item Compute for $n=3$: All three get TT then H = $\left(\frac{1}{2}\right)^9 = \frac{1}{512}$
        \item Notice pattern: contributions are $\frac{1}{8}, \frac{1}{8^2}, \frac{1}{8^3}, \ldots$
    \end{itemize}
    \item \textbf{Penultimate Step}: The final step will be summing an infinite geometric series with first term $a = \frac{1}{8}$ and ratio $r = \frac{1}{8}$
    \item \textbf{Wishful Thinking}: "I wish this were a simple geometric series..." - and it is!
    \item \textbf{Make It Easier}: First consider just two players, then extend to three
    \item \textbf{Conjecture via Small Cases}: The pattern $\frac{1}{8} + \frac{1}{64} + \frac{1}{512} + \cdots$ strongly suggests answer (B) $\frac{1}{7}$ using series formula
    \item \textbf{Visualization}: Draw a probability tree for the first few flips to see the structure
\end{itemize}

\subsection*{C. Late-Band Strategic Playbook (Problems 17-25)}
\begin{itemize}[leftmargin=*]
    \item \textbf{Answer-Choice Leverage}: Answer (A) $\frac{1}{8}$ is the probability of success on the first flip only - recognize this as a trap
    \item \textbf{Make It Easier}: Consider the two-player version first: probability would be $\frac{1}{4} + \frac{1}{16} + \frac{1}{64} + \cdots = \frac{1/4}{1-1/4} = \frac{1}{3}$
    \item \textbf{Penultimate Step}: Once we identify the series, the last step is applying the geometric series formula $\frac{a}{1-r}$
    \item \textbf{Wishful Thinking}: Assume the answer is a simple fraction, which guides us toward recognizing the geometric series structure
    \item \textbf{Conjecture via Small Cases}: Computing the first 2-3 terms reveals the pattern and confirms convergence
\end{itemize}

\subsection*{D. Argument Construction Strategy}
\begin{itemize}[leftmargin=*]
    \item \textbf{Direct Proof}: Yes - directly compute the infinite sum
    \item \textbf{Proof by Cases}: Split by number of flips $n = 1, 2, 3, \ldots$ and sum
    \item \textbf{Recursive/Inductive Approach}: Use self-similarity: $P = \frac{1}{8} + \frac{1}{8}P$ (elegant alternative)
    \item \textbf{Proof by Contradiction}: Not applicable here
    \item \textbf{Constructive Proof}: Construct the series explicitly and evaluate
\end{itemize}

\subsection*{E. Cross-Domain Translation}
\begin{itemize}[leftmargin=*]
    \item \textbf{From Probability to Algebra}: Translate the problem into summing an infinite geometric series
    \item \textbf{From Discrete to Continuous}: Not applicable (discrete is natural here)
    \item \textbf{From Combinatorics to Probability}: The problem is naturally probabilistic
    \item \textbf{From Geometry to Algebra}: Not applicable
\end{itemize}
\end{strategicbox}

\begin{tacticalbox}
\subsection*{A. Core Mathematical Tactics}
\begin{itemize}[leftmargin=*]
    \item \textbf{Primary Tactics}:
    \begin{enumerate}
        \item \textbf{Infinite Geometric Series}: This is the central technique - recognizing and summing $\sum_{n=1}^{\infty} \left(\frac{1}{8}\right)^n = \frac{1/8}{1-1/8} = \frac{1}{7}$
        \item \textbf{Independence of Events}: Use multiplication rule for independent events across three players
        \item \textbf{Law of Total Probability}: Sum over all mutually exclusive cases (different values of $n$)
    \end{enumerate}
    
    \item \textbf{Supporting Tactics}:
    \begin{itemize}
        \item \textbf{Pattern Recognition}: Identifying the geometric progression structure
        \item \textbf{Symmetry}: All three players are identical, so we can use symmetry in our reasoning
        \item \textbf{Self-Similarity/Recursion}: Alternative approach using $P = \frac{1}{8} + \frac{7}{8}P$
    \end{itemize}
\end{itemize}

\subsection*{B. Crossover Tactics}
\begin{itemize}[leftmargin=*]
    \item \textbf{Generating Functions}: Could use probability generating functions (advanced, not needed)
    \item \textbf{Markov Chains}: Could model as a Markov chain (overkill for this problem)
    \item \textbf{Recursive Equations}: The self-similarity approach creates a simple linear equation
\end{itemize}

\subsection*{C. Domain-Specific Techniques}
\begin{itemize}[leftmargin=*]
    \item \textbf{Geometric Distribution}: Probability of first success on trial $n$ is $(1-p)^{n-1}p = \left(\frac{1}{2}\right)^n$
    \item \textbf{Product of Independent Probabilities}: $P(\text{all three on trial } n) = \left[\left(\frac{1}{2}\right)^n\right]^3 = \left(\frac{1}{8}\right)^n$
    \item \textbf{Convergence of Infinite Series}: The series converges because $|r| = \frac{1}{8} < 1$
\end{itemize}
\end{tacticalbox}

\begin{conversionbox}
\subsection*{A. Recognition Index - Applicable Techniques}

\textbf{High-Probability Techniques:}
\begin{enumerate}
    \item \textbf{Infinite Geometric Series} (95\% match):
    \begin{itemize}
        \item \textit{Why it applies}: The problem naturally leads to summing $\sum_{n=1}^{\infty} ar^{n-1}$ where $a = \frac{1}{8}$ and $r = \frac{1}{8}$
        \item \textit{Recognition cue}: "Repeatedly until first success" signals geometric distribution
        \item \textit{Formula}: $\frac{a}{1-r} = \frac{1/8}{1-1/8} = \frac{1/8}{7/8} = \frac{1}{7}$
    \end{itemize}
    
    \item \textbf{Independence and Multiplication Rule} (90\% match):
    \begin{itemize}
        \item \textit{Why it applies}: Three independent coin-flipping processes
        \item \textit{Recognition cue}: "Each flips...independently"
        \item \textit{Application}: $P(A \cap B \cap C) = P(A) \cdot P(B) \cdot P(C)$ when independent
    \end{itemize}
    
    \item \textbf{Law of Total Probability} (85\% match):
    \begin{itemize}
        \item \textit{Why it applies}: Partition sample space by number of flips $n$
        \item \textit{Recognition cue}: "What is the probability..." over all possible scenarios
        \item \textit{Application}: $P = \sum_{n=1}^{\infty} P(\text{all succeed on flip } n)$
    \end{itemize}
\end{enumerate}

\textbf{Medium-Probability Techniques:}
\begin{enumerate}[resume]
    \item \textbf{Self-Similarity and Recursion} (70\% match):
    \begin{itemize}
        \item \textit{Why it applies}: After first flip, problem resets with probability $\frac{1}{8}$
        \item \textit{Alternative solution}: $P = P(\text{all heads on flip 1}) + P(\text{all tails on flip 1}) \times P(\text{same from reset})$
        \item \textit{Equation}: $P = \frac{1}{8} + \frac{1}{8}P \Rightarrow P = \frac{1}{7}$
    \end{itemize}
\end{enumerate}

\subsection*{B. Technique Selection and Integration}

\textbf{Primary Technique}: Infinite Geometric Series
\begin{itemize}
    \item This is the most direct path to the solution
    \item Requires recognizing the pattern in small cases
    \item Natural fit for the problem structure
\end{itemize}

\textbf{Secondary Technique}: Self-Similarity (Recursive Approach)
\begin{itemize}
    \item Elegant alternative that requires less computation
    \item Provides a verification method
    \item Demonstrates deeper structural understanding
\end{itemize}

\textbf{Integration Strategy}:
\begin{enumerate}
    \item Use geometric series as primary method
    \item Use recursion method as verification
    \item Both should yield $\boxed{\frac{1}{7}}$
\end{enumerate}

\subsection*{C. Technique Conflict Resolution}
\begin{itemize}[leftmargin=*]
    \item No conflicts - both approaches are complementary
    \item Geometric series is more computational, recursion is more conceptual
    \item Choose based on time constraints and comfort level
\end{itemize}
\end{conversionbox}

\begin{verificationbox}
\subsection*{A. Verification Level Selection}

\textbf{Problem Characteristics}:
\begin{itemize}
    \item Difficulty: Late-band (P19)
    \item Time pressure: Moderate (AMC12 context)
    \item Initial confidence: High after recognizing series structure
    \item Recommended level: \textbf{Level 4 (Enhanced Verification)}
\end{itemize}

\subsection*{B. Specialized Verification Methods}

\textbf{Primary Verification Methods}:
\begin{enumerate}
    \item \textbf{Alternative Method Verification}:
    \begin{itemize}
        \item Solve using recursive approach: $P = \frac{1}{8} + \frac{1}{8}P$
        \item Should confirm $P = \frac{1}{7}$
        \item If both methods agree, confidence $\rightarrow$ 95\%+
    \end{itemize}
    
    \item \textbf{Limit Checks}:
    \begin{itemize}
        \item Check: Is $0 < P < 1$? Yes, $\frac{1}{7} \approx 0.143$ $\checkmark$
        \item Check: Is $P < \frac{1}{8}$? No! Wait... $\frac{1}{7} > \frac{1}{8}$ $\checkmark$ (This makes sense because we sum multiple terms)
        \item Check: First term alone is $\frac{1}{8}$, so total must be $> \frac{1}{8}$ $\checkmark$
    \end{itemize}
    
    \item \textbf{Answer Choice Elimination}:
    \begin{itemize}
        \item (A) $\frac{1}{8}$: Too small - this is just first term
        \item (B) $\frac{1}{7}$: Matches our answer $\checkmark$
        \item (C) $\frac{1}{6}$: Too large - would require different ratio
        \item (D) $\frac{1}{4}$: Way too large - sum would diverge
        \item (E) $\frac{1}{3}$: Too large - this is the two-player answer
    \end{itemize}
    
    \item \textbf{Series Convergence Check}:
    \begin{itemize}
        \item Verify $|r| = \frac{1}{8} < 1$ so series converges $\checkmark$
        \item Verify formula application: $\frac{a}{1-r} = \frac{1/8}{7/8} = \frac{1}{7}$ $\checkmark$
    \end{itemize}
\end{enumerate}

\subsection*{C. Confidence Calibration}

\textbf{Initial Confidence}: 75\% (after recognizing geometric series)

\textbf{After Primary Verification}: 90\% (confirmed series sum)

\textbf{After Alternative Method}: 98\% (both methods agree)

\textbf{Final Confidence}: 98\% - Very high confidence to move forward
\end{verificationbox}

\begin{roadmapbox}
\subsection*{A. Step-by-Step Solution Plan}

\textbf{Solution Method 1: Infinite Geometric Series (Primary)}

\textbf{Step 1} (30 seconds): Understand the setup
\begin{itemize}
    \item Each player flips until first head
    \item Want probability all three stop on same flip
    \item Events are independent
\end{itemize}

\textbf{Step 2} (45 seconds): Compute probability for specific $n$
\begin{itemize}
    \item For one player to get first head on flip $n$: need $(n-1)$ tails then 1 head
    \item Probability: $\left(\frac{1}{2}\right)^{n-1} \cdot \frac{1}{2} = \left(\frac{1}{2}\right)^n$
    \item For all three: $\left[\left(\frac{1}{2}\right)^n\right]^3 = \left(\frac{1}{8}\right)^n$
\end{itemize}

\textbf{Step 3} (30 seconds): Set up the infinite sum
\begin{itemize}
    \item $P = \sum_{n=1}^{\infty} \left(\frac{1}{8}\right)^n$
    \item This is a geometric series with $a = \frac{1}{8}$, $r = \frac{1}{8}$
\end{itemize}

\textbf{Step 4} (30 seconds): Apply geometric series formula
\begin{itemize}
    \item Formula: $\sum_{n=1}^{\infty} ar^{n-1} = \frac{a}{1-r}$
    \item Here: $\sum_{n=1}^{\infty} \left(\frac{1}{8}\right)^n = \frac{1/8}{1-1/8} = \frac{1/8}{7/8} = \frac{1}{7}$
\end{itemize}

\textbf{Step 5} (15 seconds): Verify and select answer
\begin{itemize}
    \item Answer is $\boxed{\frac{1}{7}}$ which is choice $\textbf{(B)}$
\end{itemize}

\textbf{Total Time}: 2.5 minutes

\vspace{1em}

\textbf{Solution Method 2: Self-Similarity/Recursion (Alternative)}

\textbf{Step 1} (20 seconds): Set up recursive equation
\begin{itemize}
    \item Let $P$ = probability all three flip same number of times
    \item Consider first flip for all three players
\end{itemize}

\textbf{Step 2} (45 seconds): Analyze first flip outcomes
\begin{itemize}
    \item Case 1: All three get heads (probability $\frac{1}{8}$) → SUCCESS
    \item Case 2: All three get tails (probability $\frac{1}{8}$) → Problem resets, still has probability $P$
    \item Case 3: Mixed results (probability $\frac{6}{8}$) → FAILURE (impossible to sync)
\end{itemize}

\textbf{Step 3} (30 seconds): Write recursive equation
\begin{itemize}
    \item $P = \frac{1}{8}(\text{immediate success}) + \frac{1}{8} \cdot P(\text{reset and try again})$
    \item $P = \frac{1}{8} + \frac{1}{8}P$
\end{itemize}

\textbf{Step 4} (30 seconds): Solve for $P$
\begin{itemize}
    \item $P - \frac{1}{8}P = \frac{1}{8}$
    \item $\frac{7}{8}P = \frac{1}{8}$
    \item $P = \frac{1}{7}$
\end{itemize}

\textbf{Step 5} (15 seconds): Select answer
\begin{itemize}
    \item Answer is $\boxed{\frac{1}{7}}$ which is choice $\textbf{(B)}$
\end{itemize}

\textbf{Total Time}: 2.25 minutes

\subsection*{B. Checkpoints and Decision Points}

\textbf{Checkpoint 1} (After Step 2 of Method 1):
\begin{itemize}
    \item Verify: Did you correctly identify that each player has probability $\left(\frac{1}{2}\right)^n$?
    \item Decision: If confused, switch to Method 2 (recursion)
\end{itemize}

\textbf{Checkpoint 2} (After Step 3 of Method 1):
\begin{itemize}
    \item Verify: Is this a geometric series? Yes, check that $|r| < 1$
    \item Decision: If not confident with series, use Method 2
\end{itemize}

\textbf{Checkpoint 3} (After computing answer):
\begin{itemize}
    \item Verify: Does $\frac{1}{7}$ make sense? Is it between $\frac{1}{8}$ and $\frac{1}{4}$?
    \item Decision: If verification fails, recalculate
\end{itemize}

\subsection*{C. Common Pitfalls and Preventive Measures}

\textbf{Pitfall 1}: Forgetting to include all three players
\begin{itemize}
    \item \textit{Error}: Computing $\left(\frac{1}{2}\right)^n$ instead of $\left(\frac{1}{8}\right)^n$
    \item \textit{Prevention}: Remember to cube the probability for three independent events
    \item \textit{Detection}: Answer would be $\frac{1}{2}$ (not among choices)
\end{itemize}

\textbf{Pitfall 2}: Including only the first flip
\begin{itemize}
    \item \textit{Error}: Computing only $\frac{1}{8}$ and selecting answer (A)
    \item \textit{Prevention}: Recognize this is just the first term of an infinite sum
    \item \textit{Detection}: Problem says "repeatedly" not "once"
\end{itemize}

\textbf{Pitfall 3}: Incorrect geometric series formula
\begin{itemize}
    \item \textit{Error}: Using $\sum_{n=0}^{\infty}$ instead of $\sum_{n=1}^{\infty}$
    \item \textit{Prevention}: Be careful about series indexing
    \item \textit{Detection}: Answer would be $\frac{1}{7} + \frac{1}{8} = \frac{15}{56}$ (not among choices)
\end{itemize}

\textbf{Pitfall 4}: Arithmetic error in series sum
\begin{itemize}
    \item \textit{Error}: Miscalculating $\frac{1/8}{7/8}$
    \item \textit{Prevention}: Simplify step-by-step: $\frac{1/8}{7/8} = \frac{1}{8} \times \frac{8}{7} = \frac{1}{7}$
    \item \textit{Detection}: Use calculator or double-check arithmetic
\end{itemize}

\textbf{Pitfall 5}: Confusion in recursive approach
\begin{itemize}
    \item \textit{Error}: Forgetting the $\frac{6}{8}$ case leads to failure
    \item \textit{Prevention}: Systematically enumerate all $2^3 = 8$ outcomes of first flip
    \item \textit{Detection}: Equation would not balance correctly
\end{itemize}

\subsection*{D. Alternative Approaches}

\textbf{Alternative 1}: Direct computation of first few terms + pattern recognition
\begin{itemize}
    \item Compute: $\frac{1}{8} + \frac{1}{64} + \frac{1}{512} = \frac{64 + 8 + 1}{512} = \frac{73}{512}$
    \item Notice this is close to $\frac{1}{7} \approx 0.143$ and $\frac{73}{512} \approx 0.143$
    \item Conjecture the pattern and verify with series formula
\end{itemize}

\textbf{Alternative 2}: Answer choice testing
\begin{itemize}
    \item If answer is $\frac{1}{7}$, work backwards to verify series sum
    \item Check: Does $\frac{1}{8}(1 + \frac{1}{8} + \frac{1}{64} + \cdots) = \frac{1}{7}$?
    \item This requires $(1 + \frac{1}{8} + \frac{1}{64} + \cdots) = \frac{8}{7}$ $\checkmark$
\end{itemize}
\end{roadmapbox}

\begin{competitionbox}
\subsection*{A. Time Management Strategy}

\textbf{Problem Position}: P19 of 25 (late-band)
\begin{itemize}
    \item Expected difficulty: High
    \item Recommended time budget: 3-4 minutes
    \item Priority: Solve if comfortable with series, otherwise skip and return
\end{itemize}

\textbf{Time Allocation}:
\begin{itemize}
    \item Reading and setup: 30 seconds
    \item Solution execution: 2 minutes
    \item Verification: 30 seconds
    \item Buffer: 1 minute
    \item \textbf{Total}: 4 minutes maximum
\end{itemize}

\subsection*{B. Problem Prioritization}

\textbf{Attempt Conditions}:
\begin{itemize}
    \item $\checkmark$ Attempt if: You recognize geometric series pattern
    \item $\checkmark$ Attempt if: You're comfortable with infinite series
    \item $\checkmark$ Attempt if: You've solved P1-P18 or have 10+ minutes remaining
    \item $\times$ Skip if: Geometric series is unfamiliar
    \item $\times$ Skip if: Less than 5 minutes remaining in test
\end{itemize}

\textbf{Strategic Positioning}:
\begin{itemize}
    \item This is a "gatekeeper" problem - easier than P20-25 but harder than P1-18
    \item High success rate for students who recognize the pattern
    \item Good use of time if you see the solution quickly
\end{itemize}

\subsection*{C. Risk Management}

\textbf{Soft Abandon Triggers} (Consider switching after 5 minutes):
\begin{itemize}
    \item Cannot set up the infinite sum correctly after 2 attempts
    \item Arithmetic errors persist despite recalculation
    \item Confusion about which cases to count
\end{itemize}

\textbf{Hard Abandon Triggers} (Immediate switch):
\begin{itemize}
    \item Complete unfamiliarity with geometric series
    \item Less than 3 minutes remaining in exam
    \item Emotional frustration/panic setting in
\end{itemize}

\subsection*{D. Answer Format and Recording}

\textbf{AMC12 Multiple Choice Strategy}:
\begin{itemize}
    \item Bubble answer (B) carefully
    \item Double-check problem number (19) matches bubble sheet
    \item If unsure between (B) and another answer, use verification methods
    \item Do NOT guess randomly - if uncertain, use elimination
\end{itemize}

\textbf{Confidence-Based Strategy}:
\begin{itemize}
    \item 95\%+ confidence: Bubble and move on immediately
    \item 70-95\% confidence: Quick verification (30 sec) then move on
    \item 50-70\% confidence: Extended verification or try alternative method
    \item <50\% confidence: Consider skipping (leaving blank for 1.5 points vs. risk of 0)
\end{itemize}

\subsection*{E. Optimization Strategy}

\textbf{Expected Value Calculation}:
\begin{itemize}
    \item If 80\%+ confident: Answer is worth $0.8 \times 6 + 0.2 \times 0 = 4.8$ points
    \item If 50\% confident: Answer is worth $0.5 \times 6 + 0.5 \times 0 = 3$ points
    \item Blank worth: $1.5$ points guaranteed
    \item Decision: Answer if confidence $> 25\%$ (since $0.25 \times 6 = 1.5$)
\end{itemize}

\textbf{Time vs. Points Trade-off}:
\begin{itemize}
    \item This problem: 4 minutes for 6 points = 1.5 points/minute
    \item Easy problems: 1 minute for 6 points = 6 points/minute
    \item Decision: Prioritize easy problems first, return to P19 if time permits
\end{itemize}
\end{competitionbox}

\vspace{1em}

\begin{keyinsightbox}
\textbf{Key Insight}: The problem has elegant self-similar structure. After the first flip, if all three get tails, the problem completely resets with identical probability. This leads to both the infinite geometric series solution AND the recursive equation $P = \frac{1}{8} + \frac{1}{8}P$, both yielding $\boxed{\frac{1}{7}}$.
\end{keyinsightbox}

\begin{warningbox}
\textbf{Warning}: Answer choice (A) $\frac{1}{8}$ is a trap! This is only the probability that all three succeed on the \textit{first flip}. The problem asks for the probability over \textit{all possible numbers of flips}, which requires summing the infinite series. Don't stop after computing just the first term!
\end{warningbox}

\begin{tipbox}
\textbf{Pro Tip}: For late-band probability problems involving "repeatedly until success," immediately think geometric distribution and infinite series. The self-similarity approach ($P = \frac{1}{8} + \frac{1}{8}P$) is often faster and less error-prone than computing the series directly - learn both methods!
\end{tipbox}

\section*{Final Answer}

The probability that all three flip their coins the same number of times is:

\[\boxed{\frac{1}{7}}\]

which corresponds to answer choice $\textbf{(B)}$.

\section*{Confidence Assessment}

\begin{itemize}[leftmargin=*]
    \item \textbf{Confidence Level}: 98\% - Both solution methods (geometric series and recursion) independently confirm the answer
    \item \textbf{Verification Applied}: 
    \begin{itemize}
        \item Alternative method verification (both approaches yield $\frac{1}{7}$)
        \item Reasonableness check (answer between $\frac{1}{8}$ and $\frac{1}{3}$)
        \item Answer choice elimination (ruled out all other options)
    \end{itemize}
    \item \textbf{Flag for Review}: No - very high confidence, solution is clean and verified
    \item \textbf{Time Spent}: 2.5 minutes (actual) vs. 3-4 minutes (estimated) - excellent time management
\end{itemize}

\end{document}